\documentclass[11pt]{article}
\usepackage[margin=0.92in]{geometry}
\usepackage{amsmath,amssymb,amsthm}
\usepackage[colorlinks=true,linkcolor=blue,urlcolor=blue,citecolor=blue]{hyperref}
\newtheorem{theorem}{Theorem}
\newtheorem{lemma}{Lemma}
\newcommand{\Z}{\mathbb Z}\newcommand{\Q}{\mathbb Q}
\newcommand{\loc}{\Z_{(2)}}\newcommand{\val}{\nu_2}
\newcommand{\Cat}{\operatorname{Cat}}
\newcommand{\sig}{\sigma}
\newcommand{\coef}[1]{\left[#1\right]}
\setlength{\parskip}{4pt}
\setlength{\parindent}{1em}
\title{The dyadic denominator in shifted Jacobi parameters}
\author{}\date{Research candidate packet --- 28 October 2026}
\begin{document}
\maketitle\vspace{-2em}
\begin{abstract}
For the literal formal constants of Area's Definition 2.1, we prove that the maximum dyadic denominator exponent in independent shifted parameters is $2r+\nu_2((r-1)!)$. The argument combines a binary multidegree filtration for an integral Riccati recurrence with a normalized primitive in a quadratic polynomial algebra. The pure top coefficient attains the bound, using Area's published binomial-sum valuation. A closed Lean theorem accompanies the argument. This is preparation for independent final review; candidate credit is zero. No analytic phase identification or priority claim is made.
\end{abstract}

\section{The literal object and the statement}
Write $u=\alpha^2-1/4$, $v=\beta^2-1/4$ as independent indeterminates. Put
\[
 P_\ell(z)=\prod_{0\le j<\ell}(z-j(j+1)),\qquad
 C_s(\xi,\eta)=\sum_{\ell=0}^s
 \frac{P_\ell(u)P_{s-\ell}(v)}{\ell!(s-\ell)!}\xi^\ell\eta^{s-\ell}.
\]
Thus Area's Hankel coefficient is $a_\ell(\alpha)=P_\ell(u)/(2^\ell\ell!)$, and the above $C_s$ is exactly the coefficient in \cite[Eq.~(6)]{primary}. In the formal $Z^{-1}$-adic ring define
\begin{equation}\label{source}
 F=\sum_{s\ge0}C_s\prod_{j=1}^s(Z-j)^{-1},\quad
 L_n=\coef{Z^{-n}}\log F,\quad
 \kappa_r=\frac{\operatorname{Im}L_{2r-1}((1+i)/2,(1-i)/2)}{2^{2r-1}}.
\end{equation}
The empty product is $1$. Since the $s$th summand starts at order $s$, each coefficient uses only finitely many terms. Likewise $L_n$ is the coefficient of the finite sum $\sum_{k=1}^n(-1)^{k+1}(F-1)^k/k$. Imaginary parts treat $u,v$ as real formal variables.

For a nonzero rational $c$, its dyadic denominator exponent is $\max(0,-\val(c))$. The ring $\loc$ consists of rationals with odd reduced denominator and includes zero.
\begin{theorem}\label{main}
For every integer $r\ge1$, put $E_r=2r+\val((r-1)!)$. Every coefficient of $2^{E_r}\kappa_r(u,v)$ belongs to $\loc$. The coefficient $\coef{u^r}\kappa_r$ is nonzero and has valuation $-E_r$. Consequently the maximum dyadic denominator exponent is exactly $E_r$.
\end{theorem}
This answers the shifted question in \cite[Remark~3.8]{primary}. It concerns the power of $2$ in the common denominator; odd factors are allowed. The original-variable exponent $3r-1+\val((r-1)!)$ belongs to a different, already proved law \cite{companion}. The definition \eqref{source} is the scope of this paper. Agreement with an analytic phase expansion at every order is not proved here.

\newpage
\section{From the finite source to the midpoint recurrence}
All identities in this section are formal algebraic identities. Introduce $t$ and
\[
 \xi=(1+it)/2,\quad \eta=(1-i/t)/2,\quad
 D=-(1+t^2)\frac{d}{dt},\quad V=u(1+t^2)+v(1+t^{-2}).
\]
Here $2\xi\eta=\xi+\eta$. With
$\Delta=\xi(\xi-1)\partial_\xi+\eta(\eta-1)\partial_\eta$,
$Q=\xi^2\partial_\xi+\eta^2\partial_\eta$, and $\Theta_\xi=\xi\partial_\xi$, the product defining $P_\ell$ gives
\[
 \partial_\xi C_{s+1}=(u-\Theta_\xi(\Theta_\xi+1))C_s
\]
and the corresponding $\eta$ identity. Direct expansion gives
\begin{align*}
 \Delta Q={}&\xi(\xi-1)\Theta_\xi(\Theta_\xi+1)
 +\eta(\eta-1)\Theta_\eta(\Theta_\eta+1)\\
 &+\xi\eta(2\xi\eta-\xi-\eta)\partial_\xi\partial_\eta.
\end{align*}
The mixed term vanishes on this curve. Since $Q C_s=(\Delta+s)C_s$ by homogeneity and $D=2i\Delta$ there, we obtain
\[
 D^2C_s+2isDC_s+2iDC_{s+1}-VC_s=0.
\]
For $e_s=\prod_{j=1}^s(Z-j)^{-1}$, the identity $Ze_s=e_{s-1}+se_s$ and $DC_0=0$ now yield
\begin{equation}\label{ode} D^2F+2iZDF-VF=0.\end{equation}
Coefficientwise summation is finite, so no convergence or boundary condition enters.

Work integrally in $R=\Z[x,w]/(w^2-1-x^2)$, with unique normal form $A(x)+wB(x)$, and define
\[
 \delta x=-(1+x^2),\qquad \delta w=-wx.
\]
This derivation respects the relation. Over $\Q$, $t=w+x$, $t^{-1}=w-x$ identify this algebra with $\Q[t,t^{-1}]$, with $x=(t-t^{-1})/2$, $w=(t+t^{-1})/2$, and $D=2\delta$.
For independent integral parameters $a,b$ define
\begin{equation}\label{h}
 P=aw^2+bwx,\quad h_0=0,\quad h_1=P,\quad
 h_{n+1}=-\delta h_n-\sum_{i=1}^{n-1}h_i h_{n-i}\quad(n\ge1).
\end{equation}
Let $H_n$ be the actual half substitution $a=(u+v)/2$, $b=(u-v)/2$ in $h_n$. Then $H_1=V/4$. Dividing \eqref{ode} by $4$, and using $F(0)=1$ as a series in $Z^{-1}$, gives
\[
 \delta^2F+iZ\delta F-H_1F=0,\qquad
 \delta L_n=(-i)^n H_n.
\]
Indeed the finite logarithmic chain rule gives $F\delta\log F=\delta F$; differentiating it gives the Riccati equation and \eqref{h} by coefficient comparison. This proves the source-to-recurrence identity rather than assuming it.

\newpage
\section{The binary filtration and the two pairings}
Let $\sig(k)$ be the binary digit sum, with $\sig(0)=0$. For a nonzero bidegree $e=(i,j)$ set $q(e)=\sig(i)+\sig(j)-1\ge0$. Binary subadditivity and doubling give
\begin{equation}\label{weights}
 q(e+f)\le q(e)+q(f)+1,\qquad q(2e)=q(e).
\end{equation}
Consider parameter polynomials with zero constant coefficient whose $a^ib^j$ coefficient belongs to $2^{q(i,j)}R$. This additive class is preserved by $\delta$, by $2fg$, and by $f^2$.

For $2fg$, each product coefficient has at least the power of $2$ required by the first inequality in \eqref{weights}. For a square, distinct parameter monomials occur in pairs, giving $2f_e f_f$. A diagonal term is $f_e^2$, divisible by $2^{2q(e)}$ and hence by $2^{q(2e)}$. This proves square closure. Arbitrary multiplication is insufficient: $a,b$ lie in the class, while $ab$ does not.

A separate pairing is required in \eqref{h}: unequal spectral indices $i,n-i$ give $2h_i h_{n-i}$, and the unpaired middle term, if present, is a square. Starting with $P$, strong induction proves
\begin{equation}\label{filter}
 \coef{a^ib^j}h_n\ \in\ 2^{\sig(i)+\sig(j)-1}R\qquad (i+j>0).
\end{equation}
This induction pairs spectral indices; the proof of square closure pairs parameter monomials. Neither can be omitted.

There are two involutions. The map $\rho(x,w)=(x,-w)$ commutes with $\delta$ and replaces $b$ by $-b$ in $P$. Thus the $a^ib^j$ coefficient has $w$ parity $j$. The map $\tau(x,w)=(-x,-w)$ anticommutes with $\delta$; induction gives $\tau h_n=(-1)^{n+1}h_n$. Both follow coefficientwise from the recurrence in the torsion-free integral normal form.

All $h_n$ also belong coefficientwise to the ideal $(w)$. This holds for $P$, is preserved by $\delta(wf)=w(-xf+\delta f)$, and is preserved by the products in the recurrence. An element $A(x)+wB(x)$ belongs to $(w)$ exactly when $1+x^2$ divides $A(x)$ integrally: multiplication by $w$ sends $(C,D)$ to $((1+x^2)D,C)$.

Consequently, for odd $n$, every coefficient of $h_n$ has the integral normal form
\begin{equation}\label{normal}
 (1+x^2)A(x^2)+wxB(x^2),\qquad A,B\in\Z[X].
\end{equation}
For a coefficient of parameter degree $d=i+j$, its $B$ is zero if $j$ is even. If $j$ is odd, then
\[
 \sig(i)+\sig(j)-1=\sig(i)+\sig(j-1)\ge\sig(d-1).
\]
Equation \eqref{filter} therefore makes every coefficient of $B$ divisible by $2^{\sig(d-1)}$. The integral normal-form extraction preserves divisibility, since its polynomial divisor is monic. No comparable gain is needed for $A$.

\newpage
\section{Normalized primitives and the uniform bound}
The first summand of \eqref{normal} has the $\tau$-odd $\delta$-primitive
$-\sum_k A_k x^{2k+1}/(2k+1)$, whose midpoint value ($x=0,w=1$) is zero. For the other summand,
\[
 \delta(wf(x^2))=-wx\{f(x^2)+2(1+x^2)f'(x^2)\}.
\]
On polynomials of degree at most $N$, $Tf=f+2(1+X)f'$ is triangular with diagonal $2k+1$. Solve $Tf=-B$ downwards by
\[
 f_{N+1}=0,\qquad f_k=\frac{-B_k-2(k+1)f_{k+1}}{2k+1}.
\]
Only odd units of $\loc$ are inverted. If $2^s$ divides $B$ integrally, all $f_k$ lie in $2^s\loc$. The resulting primitive is $\tau$-odd, and its midpoint is $f_0$. Thus the normalized primitive of an $a^ib^j$ coefficient of odd $h_n$ has midpoint in
\begin{equation}\label{primitive}2^{\sig(d-1)}\loc,\qquad d=i+j>0.\end{equation}
For even $j$ that midpoint is exactly zero.

Uniqueness is essential. The kernel of $\delta$ over the rational midpoint algebra consists of scalar constants, as follows either by differentiation in the Laurent coordinate $t$ or by comparing leading coefficients in normal form. Its $\tau$-odd kernel is zero. On the literal source, $\tau$ conjugates $\xi,\eta$, hence $\operatorname{Im}L_n$ is $\tau$-odd. At $n=2r-1$ the identity on page 2 gives
\[
 \delta\operatorname{Im}L_n=(-1)^r H_n.
\]
Derivative equality and this normalization identify it with the transported primitives just constructed. There is no unproved integration constant or analytic bridge.

For $e=(i,j)$, $d=i+j$, the actual half substitution sends its monomial to
\[
 a^ib^j\longmapsto 2^{-d}(u+v)^i(u-v)^j.
\]
The numerator is integral. Consequently each coefficient of $\kappa_r$ is a finite sum of terms
\begin{equation}\label{term}(-1)^r z\,2^{-n-d}M_e,\qquad z\in\Z,\quad
 M_e\in2^{\sig(d-1)}\loc,\quad n=2r-1.\end{equation}
There is no degree-zero support. Induction in \eqref{h} gives $\deg h_n\le\lceil n/2\rceil$, so $1\le d\le r$. For the product step, odd/odd indices at even $n$ allow degree $n/2+1$; even/even allow $n/2$; mixed indices at odd $n$ allow $(n+1)/2$. These are the next-index bounds.

Legendre's identity is $\val((d-1)!)=d-1-\sig(d-1)$. With $k=\val((r-1)!)-\val((d-1)!)\ge0$, it gives
\[
 E_r+\sig(d-1)=n+d+k.
\]
Multiplying \eqref{term} by $2^{E_r}$ leaves $(-1)^r z\,2^k(M_e/2^{\sig(d-1)})\in\loc$. Finite sums prove the uniform bound for every coefficient, including zero coefficients.

\newpage
\section{Attainment and the credited binomial valuation}
Only odd/odd spectral products contribute to degree $r$ at index $2r-1$ in \eqref{h}. Their scalar recursion is Catalan's recursion, so
\[
 \coef{\text{degree }r}h_{2r-1}=(-1)^{r-1}\Cat_{r-1}P^r.
\]
To compute the pure $u^r$ midpoint, note that its coefficient in $H_1$ is $(1+t^2)/4$, and $\delta=D/2$. The odd $D$-primitive of $(1+t^2)^r$ is
$-\sum_{j=0}^{r-1}\binom{r-1}{j}t^{2j+1}/(2j+1)$.
Normalization and \eqref{source} therefore give the all-order identity
\begin{equation}\label{top}
 \coef{u^r}\kappa_r=\frac{\Cat_{r-1}}{2^{4r-2}}I_r,\qquad
 I_r=\sum_{j=0}^{r-1}\binom{r-1}{j}\frac1{2j+1}>0.
\end{equation}
This is the source's computational top formula, now obtained by the actual recurrence and primitive. Its use alone would give no upper bound for mixed coefficients.

We credit \cite[Theorem~1.4]{primary} for
\begin{equation}\label{known}\val(I_r)=r-1+\val(r).\end{equation}
The certificate proves this known ingredient internally. For completeness, here is its mathematical mechanism. Put $c_k=k!/(2k+1)!!$ in $\Q_2$. Its norm tends to zero, so all moments $M_i=\sum_{k\ge0}c_k\binom{k}{i}$ converge and have norm at most one. The recurrence $(2k+3)c_{k+1}=(k+1)c_k$, multiplied by binomial coefficients and summed, gives
\[
 (i+2)M_{i+2}=(i+1)M_i\quad(i\ge0).
\]
The boundary term is zero because $\binom0{i+1}=0$. Iteration yields
$\prod_{t<j}(2t+1)M_0=2^j j!M_{2j}$, whose norm is at most $2^{-j}$. The left product has norm one, hence $M_0=0$.

Let $R_j=\sum_{k\ge j}c_k$. Then $R_0=0$, $R_1=-1$, $R_2=-4/3$, and $\val(R_j)\ge\val(j!)$. Thus $R_j/j\in2\Z_2$ for $j\ge2$: use $R_2$ directly and $(j-1)!$ for $j\ge3$. Polynomial integration of $(1-z^2)^k$ and the binomial theorem give the finite transform
\[
 \frac{I_{m+1}}{2^m}=\sum_{k=0}^m(-1)^k\binom mk c_k
 =(m+1)\sum_{j=1}^{m+1}(-1)^j\binom m{j-1}\frac{R_j}{j}.
\]
The last sum is $1$ plus an even $2$-adic integer, proving \eqref{known} with $m=r-1$. This proof uses convergence in $\Q_2$, not real summation.

Kummer's digit formula and $(m+1)\Cat_m=\binom{2m}m$ give
$\val(\Cat_m)=\sig(m)-\val(m+1)$. Applying this and \eqref{known} to the nonzero coefficient \eqref{top}, then Legendre, yields
\[
 \val(\coef{u^r}\kappa_r)=\sig(r-1)+r-1-(4r-2)=-E_r.
\]
It follows that $2^{E_r-1}\coef{u^r}\kappa_r\notin\loc$. This completes Theorem~\ref{main}, including $r=1$, where $\kappa_1=(u-v)/4$.

\newpage
\section{Formal correspondence, dependencies, and scope}
The source is encoded over rational multivariate polynomials, with Gaussian pairs for the imaginary part. In \texttt{SourceDefs.lean}, \texttt{inverseFactorCoeff} expands each geometric factor, \texttt{powerCoeff} is finite convolution, and \texttt{sourceLogCoeff} is the finite logarithm. These definitions require no analytic interpretation. The final closed interface is
\begin{quote}\ttfamily
 JacobiShifted.source\_claim : JacobiShifted.SourceClaim.
\end{quote}
Here \texttt{SourceClaim} states the all-monomial odd-denominator bound at \texttt{dyadicExponent r} and its failure one exponent lower at \texttt{Finsupp.single 0 r}, for every natural $r\ge1$. The pure coefficient's nonzero exact valuation is separately proved.

The 35 mathematical modules are mapped in \texttt{PROOF.md} and \texttt{FORMAL-MAP.md}. The integral algebra, binary filtration, recurrence, parity and primitive modules implement pages 3--4. The literal-source, finite-log-chain and identification modules prove page 2. Degree, Catalan, pure-primitive and literal-coefficient modules prove the top identity. The three binomial modules formalize the known valuation on page 5. Finally, \texttt{DyadicArithmetic} and \texttt{HalfSubstitution} pay the degree-dependent denominator; the upper-bound, sharpness and claim modules close the theorem.

The pinned Lean executable is version 4.34.1. The packet records its binary hash, Mathlib revision, and all eight dependency revisions from the authoritative Lake manifest. Mathlib supplies quadratic algebra, polynomial and power-series operations, Catalan identities, Legendre/Kummer formulas, $2$-adic norm and convergence theory, and proof tactics. These classical ingredients receive no novelty credit. The whole-certificate audit enumerates actual environment module origins, including private/generated helpers and instances, and checks their actual types, safety flags and axiom closures. Only \texttt{propext}, \texttt{Classical.choice}, and \texttt{Quot.sound} are allowed. Audit elaborators are outside the mathematical module set.

This packet is an implementation and manuscript preparation, not independent acceptance. The proposed outcome is one family: the full shifted all-coefficient law. Candidate credit remains \textbf{0}; campaign counts remain 2/30 additional and 32/60 total. A distinct independent final review and current novelty adjudication remain pending. The preliminary source/ordinary-math review supported investment with a bounded literature screen; it did not establish worldwide priority. MathDB 391922 describes the older original-variable law and is not the target solved anew. No new record or analytic phase correspondence is asserted.

\begin{thebibliography}{9}
\bibitem{primary} I. Area, \emph{Dyadic denominator bounds and a transfer principle for the phase constants of Jacobi zeros}, arXiv:2608.23006v2, 5 October 2026. \url{https://arxiv.org/html/2608.23006v2}. Definition 2.1, (6)--(7), Theorem 1.4, and Remark 3.8.
\bibitem{companion} I. Area, \emph{The ultraspherical background of the Jacobi phase expansion, in closed form}, arXiv:2609.29531v2, 14 October 2026. \url{https://arxiv.org/html/2609.29531v2}. Theorems 1.5 and 7.36; Sections 10--11.
\end{thebibliography}
Both cited versions are public preprints; no journal reference or journal DOI is asserted. Exact cached source bodies and provenance hashes accompany this packet.
\end{document}
